\section{Introduction}

The aim is narrow on purpose: take one operation, GEMM, and push a hand written
kernel from a fair first pass to the point where Nsight Compute reports it as
compute bound and it lands within about ten percent of cuBLAS on this specific
card. The discipline that makes the result trustworthy is the comparison policy.
Percent of cuBLAS on the same device is the only headline used, because it cancels
the hardware out of the claim. No number in this report comes from anywhere but a
run of this code on this machine, and each is traceable to a results file and a
commit hash.

Two things follow from taking that policy seriously. The first is that a
measurement which contradicts the textbook is still a measurement: the shared
memory tiled kernel loses to the naive one on this part, and rather than being
skipped it becomes Section \ref{sec:negative}. The second is that a number I have
not measured yet does not get written down. Every table in this report that
covers a 1.1.0 family reads \emph{pending} in its performance columns, and the
generator that writes those tables refuses to invent a value to fill them.

The machine is an RTX 5070 (Blackwell GB205, compute capability 12.0, sm\_120, 48
SMs, 12 GB GDDR7), running CUDA Toolkit 13.3 inside WSL2. The ceilings used
throughout are measured, not read from the datasheet: about 579 GB/s streaming
bandwidth, 23 TFLOP/s FP32, and 65 TFLOP/s tensor core (best cuBLAS FP16). The
architecture reference for the part is the RTX Blackwell
whitepaper~\cite{blackwell}, and the instruction level references are the CUDA
programming guide~\cite{cudaguide} and the PTX ISA~\cite{ptxisa} shipped with
toolkit 13.3.
